\section{Introduction}
\label{sec:intro}

% The landscape of large language models (LLMs) has grown into a heterogeneous ecosystem, where frontier models coexist with far cheaper open-weight alternatives whose capabilities vary widely across tasks. Since no single model is best for every query and every budget, routing, i.e., deciding which model should answer each incoming query, has become a central mechanism for cost-effective LLM deployment. 
The rapid proliferation of large language models (LLMs) has created a heterogeneous ecosystem of models with widely varying costs and task-specific capabilities, ranging from frontier systems to substantially cheaper open-weight alternatives. Since no single model is optimal across all queries and budget constraints, model routing, which determines which model should handle each query, has become essential for cost-effective LLM deployment. Beyond cost efficiency, routing also matches each query to the candidate model best suited to it and adapts model choice to user-specific preferences (Figure~\ref{fig:main_arch}). 
Rich research has been devoted to this problem, from binary routers that arbitrate between a weak and a strong model \citep{ding2024hybrid, ong2024routellm} and cost-aware cascades \citep{chen2023frugalgpt, aggarwal2024automix}, to reward-guided ensembles, contrastive and graph-based routers \citep{chen2024routerdc, feng2024graphrouter}, personalized routers that adapt to individual users \citep{xie2025gmtrouter, dai2025personalizedrouter}, and agentic routers trained with reinforcement learning \citep{zhang2025router, feng2026graphplanner}.

Despite this rapid progress, the field still lacks a unified foundation on which these diverse approaches can be developed and compared, mainly due to two obstacles.
% First, existing routers are developed in distinct formalisms and released as separate codebases with incompatible interfaces, trained on different supervision and tuned against different candidate pools, so the design elements that actually drive performance cannot be isolated, and comparing two routers today effectively means comparing two entire experimental stacks. 
First, existing routers are developed under distinct formalisms, released as separate codebases with incompatible interfaces, trained with different supervision, and tuned for different candidate pools, making it difficult to isolate the design elements that actually drive performance or determine whether observed differences stem from the routers themselves or from their broader experimental stacks.
% Second, evaluating a router is fundamentally more demanding than evaluating a model, since both routing supervision and routing evaluation require knowing how every candidate performs on every query, which entails sweeping the full candidate pool over entire benchmarks and scoring each response with a task-specific metric. Recent benchmarks precompute candidate responses over fixed pools \citep{hu2024routerbench, huang2025routereval}, but they target one-shot routing alone and provide no pipeline for constructing new supervision, leaving multi-turn and personalized routing without a standardized, cost-aware evaluation.
Second, evaluating a router is fundamentally more demanding than evaluating a single model, as constructing routing supervision and enabling standardized evaluation require running every candidate model on every benchmark query and scoring each response with task-specific metrics. Existing benchmarks precompute candidate responses for fixed model pools \citep{hu2024routerbench, huang2025routereval}, but are limited to single-turn routing and provide no pipeline for generating supervision for new benchmarks or candidate pools. Consequently, multi-turn and personalized routing still lack a standardized, cost-aware evaluation framework, making routing methods difficult to compare, reuse, and transfer from offline studies to real applications.

In this paper, we introduce \textbf{\method}, a unified infrastructure for developing, evaluating, and deploying routing policies over heterogeneous LLM backends. Within this infrastructure, a router is characterized by five types of components: context encoders, model encoders, scoring functions, decision rules, and learning signals. This abstraction accommodates existing routers, which we group into three families of single-turn, multi-turn (including agentic), and personalized routing. It also reduces the effort to add a new router to implementing a routing method and a loss function, while data construction, training, inference, and evaluation apply unchanged, so switching the router, candidate pool, or training objective requires only a configuration change rather than reimplementation. \method\ includes more than 16 representative routers spanning all three families, and can expose any of them as an OpenAI-compatible server for deployment on messaging platforms via OpenClaw~\citep{openclaw2026} or through a ComfyUI-based visual interface for code-free prototyping.

\method\ further automates the construction of routing supervision and evaluation, the main obstacle to comparing routers on equal footing. Its pipeline assembles queries from established benchmarks, dispatches each query to a pool of 18 candidate models spanning a broad price range, and scores every response with task-specific metrics while recording token-level cost. Every router is then evaluated on the same queries, candidate pool, and metrics, enabling direct comparison of their quality-cost trade-offs. With this pipeline, we construct \textbf{xRouteBench}, a benchmark that spans generic LLM tasks, memory-augmented, vision (image and video), time-series, and personalized scenarios under one protocol.

% Leveraging the library and benchmark, we conduct a systematic empirical study of LLM routing across the three families under one protocol. Our analysis examines how far learned routing closes the gap to the best single model and to the oracle router, how the performance--cost frontier shifts across the three router families, and how much multi-turn routing and personalization add over well-tuned single-turn designs.

% We release the library, the benchmark, and all experimental artifacts at \href{https://github.com/ulab-uiuc/LLMRouter}{github.com/ulab-uiuc/LLMRouter}, hoping to spur systematic progress on LLM routing.

Leveraging \method\ and xRouteBench, we conduct a systematic empirical study of LLM routing across the three families under one protocol. Our study surfaces four findings:
% (i) \textbf{no single router dominates}, as the best router shifts from one task to the next and, on a given task, from one cost budget to the next, so the strongest average performance comes from consistency across scenarios rather than from winning any of them, and a cost-aware decision rule becomes a first-class design axis. (ii) \textbf{learned routing still beats the strongest fixed model}, since always calling the largest model is both the most expensive and only mediocre and is dominated by learned routes that recover on cheaper models the queries it fails. (iii)\textbf{ multi-turn routing brings no consistent gain on these largely single-hop workloads}, because decomposition and aggregation risk adding cost without adding information. iv) personalization pays, as conditioning the router on the user lifts accuracy on real human preferences well above the best user-agnostic design.
(i) \textbf{no single router dominates}, as the best router varies across tasks and cost budgets. Strong average performance therefore reflects consistency across scenarios rather than dominance in any single setting. (ii) \textbf{learned routing still outperforms the strongest fixed-model baseline}, because always selecting the largest model incurs the highest cost yet delivers only mediocre performance, whereas learned routers select smaller, cheaper models for many queries that the largest model answers incorrectly. (iii) \textbf{multi-turn routing does not consistently outperform single-turn routing}, as additional rounds of decomposition and aggregation often add cost and redundant information, and their benefit hinges on the capability of the base model that performs them. (iv) \textbf{personalization pays off, but only when user context is modeled well}, as a user-conditioned router ranks first under both the persona judge and real human preferences, yet the two settings favor different personalized designs.
% We release the library, the benchmark, hoping to spur systematic progress on LLM routing.
We also release the library and benchmark in the hope of fostering more systematic progress in LLM routing.

